\documentclass[11pt]{article}
\usepackage[margin=1.1in]{geometry}
\usepackage{amsmath,amssymb,amsthm,booktabs,tabularx,array}
\usepackage{microtype,xcolor}
\usepackage[colorlinks=true,linkcolor=blue!50!black,citecolor=blue!50!black,urlcolor=blue!50!black]{hyperref}
\newtheorem{theorem}{Theorem}[section]
\newtheorem{lemma}[theorem]{Lemma}
\newtheorem{proposition}[theorem]{Proposition}
\newtheorem{corollary}[theorem]{Corollary}
\theoremstyle{remark}
\newtheorem{remark}[theorem]{Remark}
\newcommand{\N}{\mathbb N}
\newcommand{\Z}{\mathbb Z}
\newcommand{\R}{\mathbb R}
\newcommand{\fr}[1]{\{#1\}}
\newcommand{\code}[1]{{\small\ttfamily #1}}
\setlength{\emergencystretch}{1.5em}
\title{Short Rotation Windows and an Infinite Family\\of Collatz Itinerary Exclusions}
\author{M.~Sharpe\thanks{Developed with Codex (OpenAI), building on the author's existing Collatz formalization. The Collatz conjecture remains open. No claim of priority in the mathematical literature is made.}}
\date{September 4, 2026}
\begin{document}
\maketitle
\begin{abstract}
For every integer $m\ge6$, set
$t_m=(\sqrt{m^2+4}-m)/2$ and $\alpha_m=1/(1+t_m)$.
We prove that no natural-number orbit of the shortcut Collatz map has
the mechanical parity itinerary of slope $\alpha_m$, for any real
intercept, even after an arbitrary finite orbit prefix.
These are distinct irrational slopes approaching one. The proof uses
an elementary short-window lemma: two positive bracketing errors
$0<\varepsilon\le\delta$ force a visit to $[1-\delta,1)$ among the first
$q+Q$ rotation points. Maximizing a fractional part proves this directly,
without a three-distance theorem. We combine this with determinant-one
separation, a uniform exponential size margin, and explicit integer
recurrences supplying every approximation certificate. The family theorem,
irrationality, strict monotonicity, and an explicit bound on the distance
to one are checked in Lean~4 using only standard foundational axioms.
This extends the repository's previous single-slope result; it does not
provide universal coverage of possible Collatz counterexamples.
\end{abstract}

\section{The family theorem}
Use the shortcut map
\[
T(n)=\begin{cases}n/2,&n\equiv0\pmod2,\\(3n+1)/2,&n\equiv1\pmod2,\end{cases}
\qquad n\in\N=\{0,1,2,\ldots\}.
\]
For real $\alpha,\rho$, its candidate mechanical itinerary is
\begin{equation}\label{eq:word}
w_{\alpha,\rho}(j)=\lfloor(j+1)\alpha+\rho\rfloor-\lfloor j\alpha+\rho\rfloor.
\end{equation}
For $0<\alpha<1$ these are binary letters; their limiting density of ones
is $\alpha$, by telescoping the floor differences.

\begin{theorem}\label{thm:main}
Let $m\in\N$, $m\ge6$, and define
\begin{equation}\label{eq:alpha}
t_m=\frac{\sqrt{m^2+4}-m}{2},\qquad
\alpha_m=\frac1{1+t_m}.
\end{equation}
For every $N,s\in\N$ and $\rho\in\R$, there exists $j\in\N$ such that
\[
T^{s+j}(N)\bmod2\ne w_{\alpha_m,\rho}(j).
\]
\end{theorem}
The formal statements are \code{Metallic.no\_\allowbreak{}itinerary}
and \code{Metallic.no\_\allowbreak{}eventual\_\allowbreak{}itinerary}.
The second follows by applying the first to $T^s(N)$; neither requires a
divergence hypothesis. Below we prove every approximation and size
obligation needed for this exclusion. For orientation only, rounded
values are
\[
\begin{array}{c|rrrrrr}
m&6&7&8&10&20&100\\\hline
\alpha_m&0.860380&0.877151&0.890388&0.909902&0.952494&0.990100
\end{array}
\]
The proof uses exact algebraic values. It covers the indicated slopes,
not all slopes between them.

\section{Short windows from a maximal fractional part}
Write $\fr{x}=x-\lfloor x\rfloor$. A \emph{visit} at time $n$ means
\[
1-\delta\le\fr{n\alpha+\rho}<1.
\]
The upper inequality holds automatically for a fractional part.

\begin{lemma}[Short-window coverage]\label{lem:window}
Suppose $q,Q$ are positive integers, $p,P\in\Z$, and
\begin{equation}\label{eq:bracket}
q\alpha=p+\delta,\qquad Q\alpha=P-\varepsilon,
\qquad 0<\varepsilon\le\delta.
\end{equation}
For every real $\rho$ there is a visit at some $0\le n<q+Q$.
Consequently, every integer window $[a,a+q+Q)$ contains a visit.
\end{lemma}
\begin{proof}
Suppose there is no visit among these finitely many indices. Choose $n$
maximizing $f(n)=\fr{n\alpha+\rho}$. Then $f(n)<1-\delta$.

If $n<Q$, the index $n+q$ is still below $q+Q$ and
$f(n+q)=f(n)+\delta$: adding $q\alpha=p+\delta$ changes the fractional
part by $\delta$, with no wrap since $f(n)+\delta<1$.
If $n\ge Q$, the index $n-Q$ is nonnegative and still in range, and
$f(n-Q)=f(n)+\varepsilon$, again with no wrap since
$f(n)+\varepsilon\le f(n)+\delta<1$.
Both cases strictly increase the maximal value, a contradiction.
Replace $\rho$ by $a\alpha+\rho$ for the shifted-window assertion.
\end{proof}
This argument requires neither coprimality nor a determinant condition.
It also respects the closed lower and open upper endpoints of the visit
interval. It proves exactly the coverage needed here, rather than a full
classification of rotation gaps. In Lean the two forms are
\code{rotation\_\allowbreak{}visit\_\allowbreak{}of\_\allowbreak{}neighbors}
and \code{rotation\_\allowbreak{}window\_\allowbreak{}of\_\allowbreak{}neighbors}.

\section{Separation and the Collatz level criterion}
For completeness, we recall the two previously formalized inputs used
with the new short window.

\begin{lemma}[Determinant-one separation]\label{lem:sep}
Under \eqref{eq:bracket}, suppose $Pq-pQ=1$. If $0<h<Q$ and $r\in\Z$,
then $|h\alpha-r|\ge\delta$.
\end{lemma}
\begin{proof}
Put $A=Ph-Qr$ and $B=qr-ph$. The determinant gives
\[
h=Aq+BQ,\qquad r=Ap+BP,\qquad h\alpha-r=A\delta-B\varepsilon.
\]
If $A\le0$, positivity of $h$ forces $B\ge1$, and $h<Q$ rules out
$A=0$. Thus $A\le-1$, so the error is at most $-\delta$.
If $A>0$, then $A\ge1$ and $h<Q$ forces $B\le0$, so the error is
at least $\delta$.
\end{proof}

\begin{proposition}[Existing endpoint criterion]\label{prop:level}
Let $\alpha\ge0$, $2\le3^\alpha$, $0<\delta<1$ and $Q\ge3$.
Suppose $q\alpha=p+\delta$, the separation conclusion of
Lemma~\ref{lem:sep} holds, and every window of length $G$ contains a
visit. If
\begin{equation}\label{eq:size}
3^{(q+G)\alpha+1}(3N+3G+1)<2^{Q-3+G},
\end{equation}
then $N$ does not have itinerary \eqref{eq:word}.
\end{proposition}
This is \code{sturmian\_\allowbreak{}level\_\allowbreak{}endpoint}, proved in
the earlier development \cite{silver,rung1}. Its proof combines the
exact affine iterate formula, including its positive correction, with
the divisibility forced by a repeated parity prefix.
Here is the mechanism needed to interpret its hypotheses. A visit
followed by a gap of at least $Q$ produces a $q$-periodic block of
$q+Q-2$ letters. Distinct orbit values with matching parity prefixes
must be far apart in ordinary integer size. The mechanical-word density
bound supplies an upper estimate for those same orbit values. The
endpoint reduction, valid when $2\le3^\alpha$, lets one check only the
largest possible starting time $G$. Equation~\eqref{eq:size} makes the
two estimates incompatible. No unproved coverage assumption will remain
when we apply this proposition.

\section{A uniform family criterion}
Call the data in \eqref{eq:bracket} a \emph{certified level} when
$q,Q>0$, $0<\varepsilon\le\delta<1$, and $Pq-pQ=1$.
These are the fields of \code{SturmianNeighborLevel}.

\begin{theorem}\label{thm:criterion}
Suppose $0\le\alpha\le1$ and $2\le3^\alpha$. Assume there is a fixed
$C\in\N$ such that, for arbitrarily large lower denominators $q$,
there are certified levels satisfying
\begin{equation}\label{eq:ratio}
6q\le Q\le Cq.
\end{equation}
Then no natural number has mechanical itinerary of slope $\alpha$,
for any real intercept.
\end{theorem}
\begin{proof}
Write $a=3^\alpha$, so $2\le a\le3$, and use the short window $G=q+Q$.
Put $d=Q-6q\ge0$. Then
\begin{equation}\label{eq:exponents}
q+G=8q+d,\qquad Q+G=13q+2d,\qquad G\le(C+1)q.
\end{equation}
Since $q>0$, we also have $Q\ge6$, so the subtraction in
\eqref{eq:size} is ordinary integer subtraction. Multiplying that
inequality by eight shows that it is equivalent to
\begin{equation}\label{eq:size24}
24(3N+3G+1)a^{q+G}<2^{Q+G}.
\end{equation}
The constants
\[
a^8\le3^8=6561<8192=2^{13},\qquad a\le4
\]
give
\begin{align*}
24(3N+3G+1)a^{q+G}
&\le24(3N+3(C+1)q+1)6561^q4^d,\\
2^{Q+G}&=8192^q4^d.
\end{align*}
Thus it suffices that
\begin{equation}\label{eq:scalar}
24(3N+3(C+1)q+1)\left(\frac{6561}{8192}\right)^q<1.
\end{equation}
For fixed $N,C$ the left side tends to zero: a geometric sequence with
base strictly below one dominates a linear factor. Choose a certified
level beyond this threshold. Lemmas~\ref{lem:window} and \ref{lem:sep}
then supply the coverage and separation hypotheses of
Proposition~\ref{prop:level}, and \eqref{eq:scalar} supplies its size bound.
All these choices are independent of the intercept.
\end{proof}
The constants are conservative. This theorem is an explicit sufficient
criterion, not a claim that arbitrary slopes have such levels. Its Lean
name is \code{no\_\allowbreak{}itinerary\_\allowbreak{}of\_\allowbreak{}unbounded\_\allowbreak{}neighbors}.

\section{The integer recurrence discharges every certificate}
Fix $m\ge6$, and abbreviate $t=t_m$, $\alpha=\alpha_m$. Direct algebra gives
\begin{equation}\label{eq:root}
t>0,\qquad t^2+mt=1,\qquad 0<t\le\tfrac12,
\qquad (1+t)\alpha=1.
\end{equation}
Hence $2/3\le\alpha<1$. In particular $3^\alpha\ge2$, since
$3^{3\alpha}\ge9>8$.

Starting from $(q_0,Q_0,p_0,P_0)=(1,1,0,1)$, define
\begin{equation}\label{eq:recurrence}
\begin{aligned}
q_{k+1}&=q_k+mQ_k,&p_{k+1}&=p_k+mP_k,\\
Q_{k+1}&=mq_{k+1}+Q_k,&P_{k+1}&=mp_{k+1}+P_k.
\end{aligned}
\end{equation}
Define the errors $\delta_k=(t^2)^k\alpha$ and
$\varepsilon_k=t\delta_k$.

\begin{lemma}\label{lem:recurrence}
For every $k\ge0$,
\begin{align}
q_k\alpha&=p_k+\delta_k,&
Q_k\alpha&=P_k-\varepsilon_k,\label{eq:errors}\\
P_kq_k-p_kQ_k&=1,&
0<\varepsilon_k&\le\delta_k<1.\label{eq:invariants}
\end{align}
Also $k+1\le q_k$, $Q_k\ge1$, and
\begin{equation}\label{eq:ratios}
Q_k\le(m+1)q_k\quad(k\ge0),\qquad
mq_k\le Q_k\quad(k\ge1).
\end{equation}
\end{lemma}
\begin{proof}
At level zero the error identities say $\alpha=\alpha$ and
$\alpha=1-t\alpha$, the latter being \eqref{eq:root}. Under one
recurrence step the lower error changes from $\delta$ to
\[
\delta-m\varepsilon=(1-mt)\delta=t^2\delta,
\]
and the upper error changes from $\varepsilon$ to
\[
\varepsilon-m(t^2\delta)=t(1-mt)\delta=t^3\delta.
\]
These are exactly the errors at the next level. Expanding
$P'q'-p'Q'$ in \eqref{eq:recurrence} gives $Pq-pQ$, preserving the
initial determinant one. Since $0<t<1$ and $0<\alpha<1$, the error
bounds follow.

The denominators stay positive and $q_{k+1}\ge q_k+1$, giving
$k+1\le q_k$. The identity $Q_{k+1}=mq_{k+1}+Q_k$ gives the lower ratio
bound. For the upper one, $m\ge1$ implies
$Q_k\le q_k+mQ_k=q_{k+1}$, and hence
$Q_{k+1}\le(m+1)q_{k+1}$; level zero is immediate.
\end{proof}

For $m=6$ the first levels are shown below. Level zero initializes the
recurrence; the sixfold ratio condition is used from level one onward.
\begin{center}
\begin{tabular}{rrrrrr}
\toprule
$k$&$q_k$&$p_k$&$Q_k$&$P_k$&$G_k=q_k+Q_k$\\
\midrule
0&1&0&1&1&2\\
1&7&6&43&37&50\\
2&265&228&1633&1405&1898\\
3&10063&8658&62011&53353&72074\\
\bottomrule
\end{tabular}
\end{center}

\begin{proof}[Proof of Theorem~\ref{thm:main}]
For the fixed parameter $m$, choose $C=m+1$. Lemma~\ref{lem:recurrence}
provides certified levels with unbounded $q_k$ and
$6q_k\le mq_k\le Q_k\le Cq_k$ for every $k\ge1$.
Theorem~\ref{thm:criterion} excludes the itinerary from every natural seed
and every intercept. Applying it to $T^s(N)$ excludes all finite-prefix
shifts as well.
\end{proof}
No continued-fraction theorem is needed to construct or validate these
levels. The induction and the positive-root equation are the complete
approximation certificates.

\section{Distinct irrational slopes approaching one}
\begin{proposition}\label{prop:geometry}
For $m\ge2$, $\alpha_m$ is irrational. The sequence $(\alpha_m)_{m\ge0}$
is strictly increasing, and for every $m\ge0$,
\begin{equation}\label{eq:close}
0<1-\alpha_m<\frac1{m+1}.
\end{equation}
In particular, the excluded slopes form an infinite set with limit one.
\end{proposition}
\begin{proof}
For $m\ge2$ the integer $m^2+4$ lies strictly between $m^2$ and
$(m+1)^2$, so it is not a square. The square root of a nonsquare
nonnegative integer is irrational; subtracting $m$, dividing by two,
adding one, and taking the reciprocal preserve irrationality here.
This uses mathlib's proved nonsquare criterion, not an asserted axiom.

If $m<n$ and $t_m\le t_n$, positivity gives
$t_m^2+mt_m<t_n^2+nt_n$, contradicting that both sides equal one.
Thus $t_n<t_m$ and $\alpha_n>\alpha_m$.

Finally $mt_m=1-t_m^2<1$, while
$1-\alpha_m=t_m\alpha_m>0$. Therefore
\[
(m+1)(1-\alpha_m)=(m+1)t_m\alpha_m
<(1+t_m)\alpha_m=1,
\]
which proves \eqref{eq:close}. The same bound, with an arbitrarily large
natural parameter, puts an excluded slope within every positive distance
of one.
\end{proof}

\section{Formal verification and scope}
The three new source modules are:
\begin{center}
\begin{tabularx}{\textwidth}{@{}lX@{}}
\toprule
Module under \code{lean/Collatz/}&Role\\
\midrule
\code{SturmianWindow.lean}&Short-window coverage for every intercept.\\
\code{SturmianFamily.lean}&Certified-level interface and uniform size criterion.\\
\code{SturmianMetallic.lean}&Recurrence certificates, family exclusion, and slope geometry.\\
\bottomrule
\end{tabularx}
\end{center}
The principal quantified type is
\begin{quote}\small\ttfamily
Collatz.Metallic.no\_\allowbreak{}itinerary (m : $\N$) (hm : 6 $\le$ m)\par
\quad (N : $\N$) ($\rho$ : $\R$) : $\neg$ Collatz.HasItin (Collatz.Metallic.alpha m) $\rho$ N.
\end{quote}
\code{HasItin} asserts equality of the parity and mechanical letters
at every natural time. The only assumptions in this type are the
parameter range and the quantified seed and intercept; there is no
remaining approximation, coverage, or asymptotic hypothesis.

The selected modules build with Lean \code{v4.27.0-rc1}.
The reproducible command
\code{python3 analysis/audit\_\allowbreak{}theorems.py --build}
checks the selected modules and prints exact types and transitive axiom
dependencies. The expanded audit contains 26 declarations, including the
new coverage, family, and geometry results. All inspected dependencies
are within \code{propext}, \code{Classical.choice}, and \code{Quot.sound}.
The new modules contain no \code{sorry}, \code{native\_\allowbreak{}decide},
or added mathematical axiom. This is ordinary Lean checking and axiom
inspection; it is not independent kernel replay, a fresh build of every
dependency, or a rebuild of every nested repository and large KL certificate.

\section{What this leaves open}
The maximal-fractional-part lemma removes a specific rotation-coverage
dependency, and the recurrence extends a single excluded slope to an
infinite explicit family. The family does not cover an interval of slopes.
Its accumulation at one gives no continuity principle that would exclude
neighboring slopes. Parameters $3,4,5$ lie outside the conservative
sixfold criterion; the earlier silver result treats the slope associated
with $m=2$ separately.

More fundamentally, no result here forces every positive-integer
counterexample to have a mechanical itinerary. General nonmechanical
itineraries, divergence at the critical boundary, and all nontrivial
cycles remain outside this family exclusion. The next proof target is a
useful extension from exact repetition to controlled weighted defects,
together with an independent restriction on possible integer
counterexample itineraries. Further exclusions alone do not supply that
coverage bridge.

\begin{thebibliography}{9}
\raggedright
\bibitem{silver} M.~Sharpe, \emph{An Unconditional Sturmian Exclusion for
the Collatz Map at Slope $1/\sqrt2$}, September 2026.
Repository file \code{paper/silver.tex}; formal sources
\code{SturmianApprox.lean}, \code{SturmianEndpoint.lean}, and
\code{SturmianSilver.lean}.
\bibitem{rung1} M.~Sharpe, \emph{Rung One of the Word-Complexity Ladder:
a Machine-Verified Exclusion of an Automatic Collatz Itinerary by
Mahler's Method in Degree One}, August 2026.
Repository file \code{paper/rung1.tex}; formal sources
\code{Rung1.lean} and \code{Sturmian.lean}.
\bibitem{source} M.~Sharpe, \emph{Collatz at the Critical Line}, source
repository, \url{https://github.com/msharpe248/collatz}.
Current research ledger: \code{analysis/RESEARCH\_\allowbreak{}STATUS.md}.
Verification metadata: \code{analysis/theorem\_\allowbreak{}audit.json}.
\end{thebibliography}
\end{document}
